% sections/00_abstract.tex
% 摘要（~200字）

\begin{abstract}
\noindent 2025年底，马斯克提出每年向轨道部署300--500\,GW AI算力，黄仁勋则认为太空散热约束将使这一愿景"需要数年"。本文以双方公开陈述为锚点，构建三层分析框架：科学层以斯特藩-玻尔兹曼定律验证1400\,\unitWpm{}散热密度的物理可行性（AI1卫星实测证实该工作点成立，但1\,GW规模需$7.1 \times 10^5$\,m\textsuperscript{2}散热面）；工程层以SpaceX/NVIDIA/NASA公开数据交叉校验发射成本、光伏效率与GPU在轨可靠性（GPU失效率是地面的3--8倍，在轨维护缺失）；商业层以TCO全口径模型对比轨道（\$681M）与地面（\$65M）10年总拥有成本，修正后的差值约为\textbf{10.5倍}。结论分三层：物理可行、工程脆弱、商业不可行。即便最乐观的2030年成本推演，轨道TCO仍为地面的4.3倍，翻转为经济可行性需发射成本、GPU能效、散热密度、硬件寿命四维度同时突破临界值，概率低于25\%。"黄仁勋的务实保留"比"马斯克的文明必然"更贴近当前技术事实。

\vspace{1em}
\textbf{研究方法}：基于斯特藩-玻尔兹曼定律的散热物理约束$\rightarrow$TCO全口径模型$\rightarrow$Wright学习率外推，三层递进验证。
\end{abstract}
